\section{Clipping cumulative segment sums}\label{clip:section}

We now perturb a primal head test while keeping the squared cost
quadratic in the perturbation size. The construction works on
$T_\infty$ and on the finite-height forest. In this section $\mathcal T$
denotes either tree system, $J(\mathcal T)$ its segment completion,
and $Z$ its closed coordinate span in $J(\mathcal T)^*$. Thus
$Z=B_\infty$ in the first case and $Z=X$ in the second. We retain
$\rho$ for the test norm and $\|\cdot\|$ for the coefficient norm.
All tests constructed below are finitely supported.

Fix a finite ancestral head $H$. For a gate $g$ let $U_g$ be its
exterior cone, and give its parent the artificial cumulative value zero
when that parent is absent from the component. For $v=Q_Hv$ define
\[
 V_g(w)=v([g,w]),\qquad
 m_g=\max_{w\in U_g}|V_g(w)|,\qquad
 q_g=\rho(v\mathbf1_{U_g}).
\]
These maxima exist: finite support gives only finitely many distinct
cumulative values. We have
\begin{equation}\label{clip:budgets}
 \sum_gm_g^2\le\rho(v)^2=\sum_gq_g^2.
\end{equation}
The equality follows from the exterior square-sum identity. For the
inequality, choose one maximizing root prefix in each nonzero cone;
these prefixes are disjoint.

A third useful grouping collects cones with the same last head vertex.
For $p\in H$, let $C_p$ consist of all exterior vertices whose last
ancestor in $H$ is $p$. Either of
\begin{equation}\label{clip:group-budgets}
 b_p=\rho(v\mathbf1_{C_p}),\qquad
 M_p=\max_{w\in C_p}|v((p,w])|
\end{equation}
may be used as its budget, with $M_p=0$ when $C_p$ is empty. Both satisfy
$\sum_pb_p^2\le\rho(v)^2$ and $\sum_pM_p^2\le\rho(v)^2$.
The first follows by joining nearly maximizing families in the disjoint
regions; for the second choose one maximizing prefix in each region.
A region may contain several cones, but a segment meeting it stays
in one cone, and disjoint crossing segments have distinct last head
vertices. These facts, rather than an assumed orthogonality within a
region, justify the grouping. For a component disjoint from $H$, there is no
last head vertex. Keep each such cone as a separate group, use its
prefix maximum or full cone norm, and choose an arbitrary fixed
inward sign. It has no head crossing and contributes only to exterior
squares. This convention also handles $H=\varnothing$.

\begin{lemma}[An inward crossing estimate]\label{clip:crossing}
Let $f=P_Hf$ and $v=Q_Hv$ be finite tests. Use either the cone budgets
$m_g$ or $q_g$, or the last-head budgets in
\eqref{clip:group-budgets}; write $c_j$ for the chosen budget of a
cone or region. Fix $s>0$. At each attached group choose the shortest
head suffix ending at its attachment whose $f$-sum has magnitude at
least $sc_j$, if one exists and $c_j>0$. Let $\sigma_j$ be the
opposite sign of that sum, and otherwise choose either sign.

Suppose a finite exterior test $w$ satisfies
$\rho(w)\le\rho(v)$, and every root prefix in group $j$ has
$w$-sum in $\sigma_j[0,c_j]$. Then
\begin{equation}\label{clip:crossing-bound}
 \rho(f+sw)^2\le\rho(f)^2+3s^2\rho(v)^2.
\end{equation}
The same conclusion holds with a strictly greater suffix threshold,
provided the small-crossing case includes equality. With a non-strict
suffix threshold, equality may also be assigned to the small-crossing
case, leaving only strictly larger terms to be split.
\end{lemma}
\begin{proof}
Take a disjoint segment family. Write $A_i$ and $B_i$ for its head
$f$-sum and exterior $w$-sum. Then
$\sum_iB_i^2\le\rho(v)^2$. A positive product $A_iB_i$ occurs
only at a crossing, and different crossings use different cone indices
(or different last head vertices in the grouped construction).
If $|A_i|<sc_j$, its cross term is bounded by $2s^2c_j^2$.
Their total is at most $2s^2\rho(v)^2$.

For a remaining positive product, the head qualifies for the suffix
selection. It contains the selected suffix, whose sum $C_i$ has the
opposite sign to $A_i$ and magnitude at least $sc_j$. The suffix is
proper, so splitting it off leaves two nonempty head segments. The
increase in their squared $f$-sums is
\begin{equation}\label{clip:split-gain}
 (A_i-C_i)^2+C_i^2-A_i^2
 =2C_i^2-2A_iC_i\ge2s|A_i|c_j\ge2sA_iB_i.
\end{equation}
These cuts are simultaneous cuts in distinct original segments.
Their pieces and the other nonempty heads are a valid disjoint test
family for $f$. Thus the original head energy plus all large positive
cross terms is at most $\rho(f)^2$. Drop nonpositive cross terms,
add the small-crossing budget $2s^2\rho(v)^2$, and add the pure
exterior budget $s^2\rho(v)^2$. This proves the result.
For a strict threshold, use $|A_i|\le sc_j$ in the small case and
$|A_i|>sc_j$ in the large case; the proof is unchanged.
\end{proof}

\begin{lemma}[Two cumulative clips]\label{clip:two-clips}
For every $s>0$ and finite tests $f=P_Hf$, $v=Q_Hv$, there are
$h_1,h_2$ supported on $\operatorname{supp}v$ such that
\begin{equation}\label{clip:difference-split}
 v=h_1-h_2,\qquad \rho(h_j)\le\rho(v),\qquad
 \rho(f+sh_j)^2\le\rho(f)^2+3s^2\rho(v)^2\quad(j=1,2).
\end{equation}
Equivalently, $v=k_1+k_2$ with $k_1=h_1$, $k_2=-h_2$, and the
two admissible perturbations are $f+sk_1$ and $f-sk_2$.
\end{lemma}
\begin{proof}
Fix any one of the budgets above and choose its signs by
Lemma~\ref{clip:crossing}. For each group let
$\Phi_j(r)=\sigma_j\max\{\sigma_j r,0\}$. On an exterior cone
in that group, set the cumulative potentials of $h_1,h_2$ equal to
$\Phi_j(V_g(w))$ and $\Phi_j(-V_g(w))$, respectively, and take
successive differences down the edges from initial value zero.
Since $\Phi_j(r)-\Phi_j(-r)=r$, the difference of the two tests
is $v$. If $v(w)=0$, its cumulative values at $w$ and its parent
are equal, so both new increments vanish. Thus no new support appears,
even along infinite descendant paths.

On any segment in a cone, its $h_j$-sum is a difference of two values
of a single $1$-Lipschitz map. Its magnitude is therefore at most the
absolute $v$-sum on that same segment. The same assertion holds for a
crossing after deleting its zero head portion. In a last-head group,
both endpoints use the same grouping index, including the artificial
zero at the head predecessor. Hence $\rho(h_j)\le\rho(v)$.
Root prefixes have the chosen sign and the chosen budget bounds their
magnitudes. Lemma~\ref{clip:crossing} completes the proof.
\end{proof}

The large-crossing calculation has a second form which retains only
one of the two split pieces. If $A_iC_i<0$, the preceding head piece
has sum $A_i-C_i$, and
\begin{equation}\label{clip:shortening}
 (A_i-C_i)^2\ge A_i^2+2s|A_i|c_j
                 \ge A_i^2+2sA_iB_i.
\end{equation}
Keeping this piece for each large crossing, and keeping each other
head unchanged, is again a disjoint head family. This proves the
homogeneous sum-decomposition estimate in
Lemma~\ref{clip:two-clips} by shortening rather than by retaining both
pieces. The missing suffix is discarded only from the test family,
not from the vector being constructed.

\subsection{What the two tests recover}\label{clip:testing}

For $x=P_Hx$ and $y\in\ker P_H$, the mixed head--tail pairings
vanish. If $\rho(f),\rho(v)\le1$, $x(f)=\|x\|$, and
$y(v)\ge\theta\|y\|$ for $0<\theta<1$, the two clips yield both
\begin{align}
 \max_\pm\|x\pm ty\|&\ge
 \frac{\|x\|+st\theta\|y\|/2}{\sqrt{1+3s^2}},
 \label{clip:max-testing}\\
 \frac{\|x+ty\|+\|x-ty\|}{2}&\ge
 \frac{\|x\|+st\theta\|y\|/2}{\sqrt{1+3s^2}}.
 \label{clip:paired-testing}
\end{align}
The first uses one clip with $|y(h_j)|\ge\theta\|y\|/2$ and
chooses the corresponding endpoint sign. The second tests $x+ty$ on
$f+sh_1$ and $x-ty$ on $f+sh_2$, and uses
$y(h_1)-y(h_2)=y(v)$ before averaging. These are distinct testing
procedures, although their displayed right sides agree. Both norms
are separately at least $\|x\|$ by testing on $f$.

Finite tail vectors have exact norming tests after a finite ancestral
projection and then $Q_H$. For completed tails, finite test density
allows $\theta\uparrow1$; alternatively approximate the tail in norm
using Lemma~\ref{binf:projections}. Thus, for every $s,t>0$,
\begin{equation}\label{clip:general-head-tail}
 \max_\pm\|x\pm ty\|\ge
 \frac{\|x\|+st\|y\|/2}{\sqrt{1+3s^2}}.
\end{equation}
This preserves arbitrary head norms and arbitrary tail magnitudes.
In particular, if $\|x-P_Hx\|<\varepsilon$ at a unit center and
$y\in\ker P_H$, $\|y\|\ge1$, contractivity gives both endpoint
bounds $\|x\pm ty\|\ge1-\varepsilon$, while
\eqref{clip:general-head-tail} applied to $P_Hx$ gives
\[
 \max_\pm\|x\pm ty\|\ge
 \frac{1-\varepsilon+st/2}{\sqrt{1+3s^2}}-\varepsilon.
\]

The paired argument also gives the same lower bound for the average.

Use the averaged modulus in \eqref{eq:averaged-modulus}, with $Z$
as the chosen coordinate space. A uniform estimate proved for finite unit centers passes to general
centers by normalized finite approximation. If its excess is $E(t)$,
choosing an approximation within $E(t)/2$ retains $E(t)/2$.
If approximation errors are allowed to tend to zero after choosing
the cofinite subspace, the full $E(t)$ remains. The supremum over
subspaces is taken before the infimum over centers. For a bound first
proved for unit tails, the symmetric norm average is an even convex
function of the tail radius, hence nondecreasing at nonnegative radii;
this gives the required $\|y\|\ge1$ version.

\begin{corollary}[The averaged clipping estimate]\label{clip:paired-modulus}
For the infinite-height coordinate predual $B_\infty$ and the
finite-height forest space $X$,
\[
 \widehat\delta_Z(t)\ge\sqrt{1+t^2/12}-1\qquad(t>0).
\]
\end{corollary}
\begin{proof}
Use the paired test procedure \eqref{clip:paired-testing}, let the
tail test approach its norm, and put $s=t/6$. At a finite unit center
the average is at least
\[
 \frac{1+st/2}{\sqrt{1+3s^2}}=\sqrt{1+t^2/12}.
\]
Arbitrarily close normalized approximation of the center preserves this
bound. The construction applies to untouched forest components because
they have no crossing terms. It needs no bound on the component height.
\end{proof}

\subsection{Homogeneous consequences of the clipped tests}
\label{dir:primal:add:homogeneous}
We record two consequences of the common clipped construction that
will be useful for comparing midpoint bounds. Here \(X\) is the finite-height
forest dual, and \(H\) is a finite ancestral head.

\begin{proposition}\label{dir:primal:add:head-tail-lens}
If \(x=P_Hx\) and \(y=Q_Hy\), then
\begin{equation}\label{dir:primal:add:optimized}
 \max_{\pm}\|x\pm y\|\ge\sqrt{\|x\|^2+\|y\|^2/12}.
\end{equation}
If \(0<\sigma\le1\) and \(w,Z_1,Z_2,\ldots\in X\) satisfy
\(\|w\pm Z_j\|\le1\) and \(\|Z_j-Z_l\|\ge2\sigma\) for \(j\ne l\),
then \(\|w\|\le1-\sigma^2/24\).
\end{proposition}
\begin{proof}
In \eqref{clip:general-head-tail}, put \(t=1\) and
\(s=\|y\|/(6\|x\|)\) when both vectors are nonzero. This gives
\eqref{dir:primal:add:optimized}. If \(y=0\) the assertion is equality;
if \(x=0\), it follows from \(\|y\|\ge\|y\|/\sqrt{12}\).

For the second assertion, given \(\varepsilon>0\), choose \(H\) with
\(\|w-P_Hw\|\le\varepsilon\). The bounded sequence \(P_HZ_j\) has two
distinct terms for which \(h=(Z_j-Z_l)/2\) satisfies
\(\|P_Hh\|\le\varepsilon\). Then \(\|h\|\ge\sigma\) and
\(\|w\pm h\|\le1\) by convexity. Formula~\eqref{clip:general-head-tail}
at \(P_Hw,Q_Hh\), again with \(t=1\), gives
\[
 (1+2\varepsilon)\sqrt{1+3s^2}
 \ge\|w\|-\varepsilon+\tfrac{s}{2}(\sigma-\varepsilon).
\]
Let \(\varepsilon\downarrow0\), put \(s=\sigma/6\), and use
\(\sqrt{1+3s^2}\le1+3s^2/2\). This retains the stated \(1/24\) loss.
\end{proof}

\subsection{Three ways to retain a useful sign}\label{clip:selection}

We return to both choices of $Z$ for the three selection procedures below.

The two global clips retain half the original pairing in a single
endpoint test. Sometimes the construction selects signs before it
chooses that endpoint. The resulting extra loss records a different
selection mechanism.

\begin{proposition}[Selection before testing]\label{clip:selection-prop}
Suppose $x=P_Hx$, $\|x\|=1$, $y\in\ker P_H$, and
$v=Q_Hv$, $\rho(v)\le1$, and $y(v)>\|y\|/2$.
For every $s>0$, each of
the following constructions produces an exterior test $w$ with
\[
 |y(w)|>\|y\|/8,\qquad \rho(w)\le1,
 \qquad \rho(f+sw)^2\le1+3s^2,
\]
where $f=P_Hf$ norms $x$ and the inward signs are chosen with full
cone budgets $q_g$.

\begin{enumerate}
\item Take the positive and negative cumulative parts globally, choose
one by its total pairing, then retain one of the two groups of cones
having the same prescribed inward sign.
\item Form all four combinations of a cumulative positive/negative
part and a prescribed inward-sign group, and select one by its pairing.
\item Choose a cumulative part separately in each cone, orient it
inward, then retain the cones on which the resulting pairings have
one common sign.
\end{enumerate}
\end{proposition}
\begin{proof}
The exterior-test requirement loses no norming accuracy: replacing any
finite test $v$ by $Q_Hv$ preserves $y(v)$ and cannot increase $\rho(v)$.
Let $v^+,v^-$ be the successive differences of $\max(V_g,0)$ and
$\max(-V_g,0)$ on each cone. Both have nonnegative root prefixes,
$v=v^+-v^-$, and each segment sum is dominated in absolute value
by the corresponding $v$-sum.

For the first construction choose $q\in\{v^+,v^-\}$ with
$|y(q)|>\|y\|/4$. Partition its cones by $\sigma_g=1$ and
$\sigma_g=-1$, writing $q=q_++q_-$. One part has pairing of
absolute value greater than $\|y\|/8$. For that sign $\sigma$
put $w=\sigma q_\sigma$. This last sign makes its root prefixes
inward; the eventual endpoint sign is chosen from the sign of $y(w)$.
These two choices must not be confused.

For the second construction form
$w_{\sigma,+}=\sigma v^+\mathbf1_{\{\sigma_g=\sigma\}}$ and
$w_{\sigma,-}=\sigma v^-\mathbf1_{\{\sigma_g=\sigma\}}$ for
$\sigma\in\{-1,1\}$. Then
$v=\sum_\sigma\sigma(w_{\sigma,+}-w_{\sigma,-})$.
The triangle inequality for their four pairings gives the claimed
$1/8$ bound for one of them. Each of the four was fixed before the
endpoint was selected.

For the third construction write $v_g=v\mathbf1_{U_g}$. In each
nonzero cone select $q_g'$ from the two cumulative parts so that
$|y(q_g')|\ge|y(v_g)|/2$. Reverse the selected part if necessary
to make its root prefixes inward. Call the oriented test $r_g$.
Then
\[
 \sum_g|y(r_g)|\ge\tfrac12\sum_g|y(v_g)|
                    \ge\tfrac12y(v)>\|y\|/4.
\]
Keep either the positive-pairing cones or the negative-pairing cones,
whichever has larger total absolute pairing, and sum their $r_g$.
This loses another factor two. The selected half in each cone can
differ; this procedure is not the choice of one of two fixed global
clips.

In all three constructions, masking cones and reversing complete
cone pieces preserves the same-segment absolute-sum domination.
Thus $\rho(w)\le\rho(v)$ and all prefix signs are inward.
Lemma~\ref{clip:crossing} proves the test-cost bound.
\end{proof}

For each selection procedure the selected endpoint satisfies
\begin{equation}\label{clip:eighth-ratio}
 \max_\pm\|x\pm ty\|\ge
 \frac{1+st\|y\|/8}{\sqrt{1+3s^2}},
\end{equation}
and the other endpoint has norm at least one. If $\|y\|\ge1$, then for
every fixed $t>0$ the ratio exceeds one by a uniform amount when $s>0$
is sufficiently small. The previously
described center approximation then gives a uniform positive averaged
gain. This route spends a second factor two to impose the selected sign;
the paired procedure uses both test evaluations instead.

\subsection{Selection order and whole-crossing replacement}
\label{dir:primal:add:selection}
In this subsection we specialize to the finite-height forest dual $X$.
The constructions in Proposition~\ref{clip:selection-prop} allow two
further choices worth distinguishing. The first chooses its prefix budget
only after selecting a global clipped part. The second pays for a whole
crossing square with a single shortened head segment.

\paragraph{Choosing the budget after a global clip.}
\label{dir:primal:add:global}
Let \(0<t\le1\), \(\|x\|=1\), and put \(\eta=t^2/1000\).
Choose a finite ancestral \(H\) with \(\|Q_Hx\|<\eta\), and a head test
\(f\) with \(\rho(f)\le1\), \(x(f)>1-2\eta\).
For \(z\in\ker P_H\) with \(\|z\|\ge t\), choose a finite unit-ball
tail test \(v\) with \(z(v)>3t/4\).
Take positive and negative cumulative parts as in
Proposition~\ref{clip:selection-prop}. One global part \(q\) has
\(|z(q)|>3t/8\); all its root prefixes are nonnegative. Now assign to
each cone the prefix maximum of \(q\), rather than that of \(v\),
\[
 M_g=\max_{u\in U_g}q([g,u]),\qquad \sum_gM_g^2\le1.
\]
With \(s=t/32\), choose the inward sign from the shortest head suffix
of magnitude at least \(sM_g\), taking a default sign when \(M_g=0\),
there is no parent, or there is no qualifying suffix.
One of the two sign classes of cones retains absolute pairing greater
than \(3t/16\). Multiply that class by its common inward sign, and call
it \(w\). Then \(\rho(w)\le1\). Lemma~\ref{clip:crossing}, applied with
baseline \(q\) and its own prefix budgets, gives
\(\rho(f+sw)^2\le1+3s^2\). Since \(|x(w)|<\eta\) and \(z(f)=0\),
\begin{equation}\label{dir:primal:add:global-gain}
 \max_{\pm}\|x\pm z\|
 \ge\frac{1-2\eta-s\eta+3st/16}{\sqrt{1+3s^2}}
 \ge1+t^2/1000.
\end{equation}
Indeed, after replacing the denominator by \(1+3s^2/2\), the numerator
minus \((1+t^2/1000)(1+3s^2/2)\), divided by \(t^2\), is
\[
 \frac3{512}-\frac3{1000}-\frac t{32000}
 -\frac3{2048}-\frac{3t^2}{2048000}
 \ge\frac{2789}{2048000}>0.
\]
Thus this order of selection keeps both losses \(3t/8\) and \(3t/16\).

\paragraph{One head piece for a whole crossing.}
\label{dir:primal:add:whole-crossing}
Take a finite head test \(f\) with \(\rho(f)\le1\), a unit
\(y\in\ker P_H\), and a finite tail test \(v\) with
\(\rho(v)\le1\), \(y(v)>1/2\).
The third construction in the proof of Proposition~\ref{clip:selection-prop}
uses only \(f\) and these tests, before any center is evaluated. With full
cone budgets \(q_g\) and threshold \(s q_g\), it produces \(w\) with
\(\rho(w)\le1\), \(|y(w)|>1/8\), and inward root prefixes of magnitude
at most \(q_g\).

Consider a crossing segment with head \(P\), and put
\(A=f(P)\), \(B=w([g,u])\). If its cone is unselected or \(B=0\),
keep the head square. Otherwise \(q_g>0\). If \(|A|<s q_g\), its whole
square is at most \(A^2+3s^2q_g^2\). If \(|A|\ge s q_g\), the selected
suffix \(I\subseteq P\), with sum \(C=f(I)\), exists. When \(A\) and
\(C\) have the same sign, the inward tail cancels and
\(|A+sB|\le|A|\). With opposite signs, \(I\) is proper and
\[
 |f(P\setminus I)|=|A|+|C|\ge|A+sB|.
\]
In that case replace the whole head test by its preceding piece.
All retained head pieces are disjoint. Small crossings cost at most
\(3s^2\sum_gq_g^2\), and wholly exterior segments cost at most \(s^2\).
We have proved, directly by whole-crossing replacement,
\begin{equation}\label{dir:primal:add:four-cost}
 \rho(f+sw)^2\le1+4s^2.
\end{equation}
This differs from paying the cross term and tail square separately in
\eqref{clip:shortening}.

For an actual unit center \(x\), take
\(s=t/64\) and \(\eta=\min\{1/2,ts/[32(1+s)]\}\), for any \(t>0\).
Choose \(f\) finite with \(\rho(f)\le1\), \(x(f)>1-\eta\), and a finite
coordinate approximation \(x_0\) within \(\eta\) of \(x\). Choose \(H\)
to contain both supports and their ancestors. For each unit tail \(y\),
the selected \(w\) above satisfies \(|x(w)|<\eta\). Hence
\begin{equation}\label{dir:primal:add:four-gain}
 \max_{\pm}\|x\pm ty\|
 \ge\frac{1-(1+s)\eta+ts/8}{\sqrt{1+4s^2}}
 \ge\frac{1+3ts/32}{1+ts/32}>1.
\end{equation}
The last step uses \(2s^2=ts/32\).

\subsection{Direct estimates at an approximate center}
\label{clip:approximate-center}

Here $Z$ again denotes either coordinate model. We record the error
calculation because it permits the tests to be
chosen at the actual center, and it distinguishes paired testing from
selection of one endpoint. Let $\|x\|=1$, let $x_0=P_Hx_0$ with
$\|x-x_0\|\le\varepsilon$, and choose a head test $f$ with
$\rho(f)\le1$. Suppose
\[
 x(f)\ge1-\kappa\varepsilon,
 \qquad y\in\ker P_H,\qquad \|y\|\ge1.
\]
Every exterior unit test $w$ satisfies $|x(w)|\le\varepsilon$.
If a selected inward test has $|y(w)|\ge\theta$, then
\begin{equation}\label{clip:approx-selected}
 \max_\pm\|x\pm ty\|\ge
 \frac{1-\kappa\varepsilon-s\varepsilon+st\theta}
      {\sqrt{1+3s^2}},
 \qquad \min_\pm\|x\pm ty\|\ge1-\kappa\varepsilon.
\end{equation}
These follow by evaluating on $f+sw$ with the chosen endpoint sign,
and on $f$ itself. If $v=h_1-h_2$ and $y(v)\ge\theta$, testing
both endpoints instead gives
\begin{equation}\label{clip:approx-paired}
 \frac{\|x+ty\|+\|x-ty\|}{2}
 \ge\frac{1-\kappa\varepsilon-s\varepsilon+st\theta/2}
           {\sqrt{1+3s^2}}.
\end{equation}
Here the tail pairings combine exactly to $y(v)$, while the average
of the two center errors is bounded below by $-s\varepsilon$.
If one instead works first at $x_0$, with $x_0(f)\ge1-2\varepsilon$,
and then transfers the norm estimate back to $x$, the corresponding
paired bound is
\begin{equation}\label{clip:finite-transfer}
 \frac{\|x+ty\|+\|x-ty\|}{2}
 \ge\frac{1-2\varepsilon+st\theta/2}{\sqrt{1+3s^2}}
       -\varepsilon.
\end{equation}

\paragraph{Choosing the test size before its accuracy.}
\label{clip:ordered-accuracy}
The formulas also prove positivity without choosing a numerical test
size in advance. With a half-norming tail and one global clip, first
choose $0<s<1$ so that
$(1+st/4)/\sqrt{1+3s^2}>1$. Then choose the center and head-test
accuracy $\varepsilon$ small enough that
\[
 \frac12\left[
 \frac{1-\varepsilon-s\varepsilon+st/4}{\sqrt{1+3s^2}}
      +1-\varepsilon\right]>1.
\]
The head normer and its finite approximation are chosen to this
accuracy, and the cofinite subspace annihilates the resulting head.
The choices depend on $t$, not on the unit tail. Their order separates
the quadratic test cost from the approximation error.

\paragraph{Using a norming-functional kernel.}\label{clip:norming-kernel}
Let $\ell\in Z^*$ norm the actual center $x$, and use
$F=\ker P_H\cap\ker\ell$. Then both endpoints have norm at least
one, without a center error in that conclusion. A tail test with
pairing above $3/4$ supplies one global clip with pairing above $3/8$.
Choosing a head test with $x(f)\ge1-3\varepsilon$, formula
\eqref{clip:approx-selected} gives
\[
 \max_\pm\|x\pm ty\|\ge
 \frac{1-3\varepsilon-s\varepsilon+3st/8}{\sqrt{1+3s^2}}.
\]
Averaging this bound with one exploits information from the added
kernel that the finite-coordinate conditions alone do not supply.

\paragraph{Testing the two endpoints directly.}\label{clip:direct-paired}
For a half-norming tail, the difference decomposition and two head-plus-tail
tests instead give
\[
 \frac{\|x+ty\|+\|x-ty\|}{2}
 \ge\frac{1-3\varepsilon-s\varepsilon+st/4}{\sqrt{1+3s^2}}.
\]
Here the tail gain is the pairing of the original tail test, divided
by two; the center errors are estimated before averaging. Thus this
procedure controls the sum directly. Formula~\eqref{clip:finite-transfer}
records the other order: test a finite approximant first, then transfer
both endpoint norms to the actual center.

\paragraph{Separating crossing and wholly exterior budgets.}
\label{clip:constant-four}
For full-cone budgets, keep the sum decomposition $v=k_1+k_2$.
One may bound a crossing by its retained or split head square plus
$3s^2q_g^2$: $2s^2q_g^2$ pays the small cross term and
$s^2q_g^2$ pays its tail square. Sum over crossings and then pay
at most another $s^2$ for all wholly exterior segments. This gives
\[
 \rho(f+sk_1)^2,\ \rho(f-sk_2)^2\le1+4s^2.
\]
With a half-norming tail and a finite-center test of value at least
$1-2\eta$, paired testing followed by center transfer therefore gives
\[
 \frac{\|x+ty\|+\|x-ty\|}{2}
 \ge\frac{1-2\eta+st/4}{\sqrt{1+4s^2}}-\eta.
\]
First choose $s>0$ so that the ratio at $\eta=0$ exceeds one,
then choose $\eta$ from this positive gap. This retains the separate
crossing and exterior accounting, including its constant $4$.

